\documentclass[Main.tex]{subfiles}

%these make the chapter/section numbers correct if built individually
%obviously you'd have to adjust these to be accurate.
\setcounter{chapter}{2}
\setcounter{section}{3}

\externaldocument{Literature Review}

\begin{document}
\ourchapter{Literature Review}\label{ch:literature review}
% summarise

Words, including anthroponyms and toponyms, can be seen as sequences of characters. Traditional feedforward networks have typically shown to be unsuitable for this task for several reasons. First, they expect a fixed-size input, which is inappropriate or, at least, inconvenient when it comes to varying-length sequences, such as words. Second, words are such that the order of characters matters, so training a word as a "bag of characters" would lose all positional information. Finally, this type of networks treats characters in different parts of the word as different due to lack of weight sharing. In the perhaps more intuitive example of a sentence, the meaning of a word should generally and approximately be the same whether it appears at the beginning or at the end of a sentence. As a result, the task of sequence modelling has often become a task of temporal sequence modelling. However, traditional feed-forward networks have proven unsuitable even for this task, as they lack an explicit mechanism for representing time if not as an additional dimension of the input, which leads to similar shortcomings as the ones mentioned above relatively to the use of feed-forward networks for the more generic task of sequence modelling.\\

Recurrent Neural Networks ("RNNs"), introduced by \cite{elman_finding_1990}, represent time by incorporating recurrent connections that, at each timestep, feed a hidden state back into the network, effectively giving the network a form of memory. When applied to a sequential task such as name generation, this formulation is autoregressive: the sequence is generated one character at a time, with each character conditioned on all preceding characters via a hidden state, that carries forward a summary of the sequence seen so far. 

In a modern vanilla RNN used in the context of character sequence generation, at each timestep, the hidden state is updated as a function of the current input and the previous hidden state, with the recurrent connection acting as an implicit memory.
However, vanilla RNNs suffer from the vanishing or exploding gradient problem, whereby gradients tend to diminish or blow up exponentially as they are backpropagated through time, making it difficult to capture dependencies across longer sequences.\\

Long Short-Term Memory networks ("LSTMs"), introduced by \cite{hochreiter_long_1997}, have gained increasing popularity over vanilla RNNs thanks to their ability to learn long-term dependencies and are now widely used for a variety of tasks. Just like vanilla RNNs, they are also characterised by a repeating module, however this module contains three gating mechanisms, i.e., a "forget" gate, an "input" gate and, finally, an "output" gate. To the hidden state that is typical of RNNs, LSTMs add a cell state, which is responsible for long-term memory, leaving short-term memory a responsibility of the hidden state. At each timestep, the network updates the cell state by deciding, based on the input and on the previous hidden state, what to forget of the previous cell state and what to remember of the input. The network also updates the hidden state based on the context, represented by the output gate and the cell state.\\

As shown by Graves \cite{graves_generating_2013}, LSTMs can be used to generate sequences of characters autoregressively. At each timestep, the network takes a character embedding as input and outputs a probability distribution over the vocabulary of possible next characters, produced by passing the hidden state through a softmax activation function. However, the question arises as to what input should be fed into the network. During training, teacher forcing is often used, so the output is ignored and the network is instead given the correct character or, rather, its embedding, as input. At inference time, a common approach consists in using greedy decoding, whereby the character with the highest probability is always picked and fed back into the network. \cite{graves_generating_2013} proposes an interesting alternative to the greedy approach which consists in sampling from the network's output distribution and feeding the sample back to the network. While the network itself remains deterministic since weights and biases are fixed, this sampling step introduces stochasticity into the model. When RNNs are used as building blocks of generative models, this "fuzziness" can become a strength of the model, allowing it to produce diverse and novel outputs. Another approach consists in using beam search, whereby a number $B$ of partial sequences is kept at each step based on a score, kept at each step based on a score, trading off computational cost for higher quality sequences compared to greedy decoding.\\

Autoencoders constitute a type of neural network architecture designed to efficiently compress (or encode) input data, then reconstruct (or decode) the original input from this compressed representation, also called code or latent representation. In other words, they learn the identity function in a self-supervised way. In training, autoencoders are optimized via backpropagation to minimize reconstruction loss, i.e., the discrepancy between the reconstructed data point and the original input data. This forces the model to learn a compact and meaningful latent representation. Cross-entropy loss can be used as reconstruction loss. Autoencoders are typically composed of an encoder, a "bottleneck" and a decoder. The encoder comprises layers that achieve an increasingly compressed representation of the input data through dimensionality reduction. The bottleneck (or “code”) contains the most compact representation of the input: it is both the output layer of the encoder network and the input layer of the decoder network. The decoder, finally, decompresses the encoded representation of data, ultimately reconstructing the data back to its original, pre-encoding form. This reconstructed output is then compared to the original input to measure model performance.\\

When Recurrent Neural Networks are used as building blocks for an autoencoder, the architecture is typically referred to as a Recurrent Autoencoder (RAE) or Sequence-to-Sequence (Seq2Seq) model \cite{sutskever_sequence_2014}. In these models, both the encoder and the decoder are RNNs, and sequences are generated autoregressively, one token at a time.
The encoder reads the input sequence character by character. At the first timestep, the encoder takes a start-of-sequence token as input and initialises its hidden state to a zero vector. At each subsequent timestep, each layer of the encoder takes as input the character embedding of the current token - or, for layers above the first, the hidden state of the layer below at the same timestep — and updates its hidden state accordingly. The intermediate hidden states of the encoder are discarded; only the final hidden state of the top layer, produced after processing the last token of the input sequence, is retained. This vector, of fixed dimensionality regardless of the length of the input sequence, constitutes the bottleneck or latent representation, and is the only information passed to the decoder. The decoder is initialised with this latent vector as its hidden state. It then generates the output sequence autoregressively, one character at a time, conditioned on the latent vector and on the previously generated character. At training time, teacher forcing is typically employed, whereby the correct character is fed as input at each timestep regardless of the model's predictions. At inference time, the decoder generates characters by sampling from or greedily decoding its output distribution, until an end-of-sequence token is produced. he reconstruction loss of a sequence autoencoder is computed as the sum of cross-entropy losses across all timesteps of the output sequence. At each timestep $t$, the decoder outputs a probability distribution over the vocabulary of possible next characters via a softmax activation function. The cross-entropy loss measures the discrepancy between this predicted distribution and the true next character, represented as a one-hot vector. The total reconstruction loss using teacher forcing is therefore:
\begin{equation}
\mathcal{L}_{\text{reconstruction}} = -\sum_{t=1}^{T} \log p_\theta(y_t \mid y_{1}, \ldots, y_{t-1}, \mathbf{z})
\end{equation}
where $\mathbf{z}$
is the latent vector produced by the encoder and $y_t$ is the true character at timestep
$t$. Since the output sequence includes start and end of sequence tokens, total sequence length for a name of length $L$ would be $T=L+2$. This loss is equivalent to maximising the log-likelihood of the correct output sequence given the latent representation and given the previous ground truth characters, encouraging the decoder to faithfully reconstruct the original input.\\

%choice of token
%character embeddings

%\tocsubsectionstar{Future Work} %optional

%\subsection*{Closing Comments}

\end{document}
